% ============================================================
%  COMFHA Lab --- Command Line & HPC Cluster Tutorial
%  For new lab members | Kasetsart University
% ============================================================
\documentclass[12pt,a4paper]{article}

% -- Packages ------------------------------------------------
\usepackage[top=2.5cm,bottom=2.5cm,left=3cm,right=3cm]{geometry}
\usepackage[T1]{fontenc}
\usepackage[utf8]{inputenc}
\usepackage{lmodern}
\usepackage{microtype}
\usepackage{xcolor}
\usepackage{mdframed}
\usepackage{listings}
\usepackage{hyperref}
\usepackage{titlesec}
\usepackage{titletoc}
\usepackage{fancyhdr}
\usepackage{graphicx}
\usepackage{booktabs}
\usepackage{array}
\usepackage{enumitem}
\usepackage{amsmath}
\usepackage{tcolorbox}
\tcbuselibrary{skins,breakable,listings}
\usepackage{parskip}
\usepackage{setspace}

% -- Colors --------------------------------------------------
\definecolor{kugreen}{RGB}{0,104,56}
\definecolor{kugold}{RGB}{187,148,0}
\definecolor{codebg}{RGB}{245,247,250}
\definecolor{codeframe}{RGB}{180,195,215}
\definecolor{tipbg}{RGB}{232,247,235}
\definecolor{tipframe}{RGB}{0,140,80}
\definecolor{warnbg}{RGB}{255,248,220}
\definecolor{warnframe}{RGB}{200,150,0}
\definecolor{terminalBG}{RGB}{30,30,30}
\definecolor{terminalFG}{RGB}{204,204,204}
\definecolor{promptgreen}{RGB}{80,200,120}
\definecolor{commentgray}{RGB}{130,130,130}
\definecolor{keywordblue}{RGB}{86,156,214}
\definecolor{stringorange}{RGB}{206,145,120}

% -- Hyperref ------------------------------------------------
\hypersetup{
  colorlinks = true,
  linkcolor  = kugreen,
  urlcolor   = kugreen,
  citecolor  = kugreen,
  pdfauthor  = {COMFHA Lab, Kasetsart University},
  pdftitle   = {Command Line and HPC Cluster Tutorial},
}

% -- Section styling -----------------------------------------
\titleformat{\section}
  {\large\bfseries\color{kugreen}}
  {\thesection.}{0.8em}{}[\titlerule]

\titleformat{\subsection}
  {\normalsize\bfseries\color{kugreen!80!black}}
  {\thesubsection.}{0.6em}{}

\titleformat{\subsubsection}
  {\normalsize\bfseries}
  {\thesubsubsection.}{0.5em}{}

% -- Header / Footer -----------------------------------------
\pagestyle{fancy}
\fancyhf{}
\renewcommand{\headrulewidth}{0.4pt}
\renewcommand{\footrulewidth}{0.4pt}
\fancyhead[L]{\small\color{kugreen}\textbf{COMFHA Lab} --- Command Line \& HPC Tutorial}
\fancyhead[R]{\small\color{gray}Kasetsart University}
\fancyfoot[C]{\small\thepage}
\fancyfoot[R]{\small\color{gray}For internal use}

% -- Terminal box (dark theme) --------------------------------
\lstdefinestyle{terminal}{
  backgroundcolor  = \color{terminalBG},
  basicstyle       = \ttfamily\small\color{terminalFG},
  breaklines       = true,
  breakatwhitespace= true,
  frame            = single,
  rulecolor        = \color{terminalBG!70!black},
  framesep         = 6pt,
  xleftmargin      = 8pt,
  xrightmargin     = 8pt,
  aboveskip        = 8pt,
  belowskip        = 8pt,
  showstringspaces = false,
  moredelim        = [is][\color{promptgreen}\bfseries]{<G>}{</G>},
  moredelim        = [is][\color{commentgray}\itshape]{<C>}{</C>},
  moredelim        = [is][\color{stringorange}]{<O>}{</O>},
  moredelim        = [is][\color{keywordblue}]{<B>}{</B>},
}

\lstdefinestyle{code}{
  backgroundcolor  = \color{codebg},
  basicstyle       = \ttfamily\small,
  breaklines       = true,
  frame            = single,
  rulecolor        = \color{codeframe},
  framesep         = 5pt,
  xleftmargin      = 8pt,
  xrightmargin     = 8pt,
  aboveskip        = 6pt,
  belowskip        = 6pt,
  showstringspaces = false,
}

\lstdefinestyle{bash}{
  style            = code,
  morekeywords     = {ssh,scp,rsync,sbatch,squeue,scancel,sinfo,module,source,
                      conda,ls,cd,pwd,mkdir,rm,cp,mv,cat,less,head,tail,nano,
                      chmod,ssh-keygen,ssh-copy-id,grep,find,echo,export},
  keywordstyle     = \color{kugreen!80!black}\bfseries,
  commentstyle     = \color{commentgray}\itshape,
  comment          = [l]{\#},
}

% -- Tip / Warning boxes -------------------------------------
\tcbset{
  tipstyle/.style={
    enhanced, breakable,
    colback=tipbg, colframe=tipframe,
    fonttitle=\bfseries\small,
    title={Tip},
    left=6pt, right=6pt, top=4pt, bottom=4pt,
    before upper={\setlength\parskip{4pt}},
  },
  warnstyle/.style={
    enhanced, breakable,
    colback=warnbg, colframe=warnframe,
    fonttitle=\bfseries\small\color{warnframe!80!black},
    title={$\bigtriangleup$\ Important},
    left=6pt, right=6pt, top=4pt, bottom=4pt,
    before upper={\setlength\parskip{4pt}},
  },
  exercisestyle/.style={
    enhanced, breakable,
    colback=white, colframe=kugreen,
    fonttitle=\bfseries\small\color{kugreen},
    title={Exercise},
    left=6pt, right=6pt, top=4pt, bottom=4pt,
    before upper={\setlength\parskip{4pt}},
  },
  defstyle/.style={
    enhanced, breakable,
    colback=blue!5, colframe=blue!40,
    fonttitle=\bfseries\small,
    left=6pt, right=6pt, top=4pt, bottom=4pt,
    before upper={\setlength\parskip{4pt}},
  },
}

\newcommand{\tip}[1]{%
  \begin{tcolorbox}[tipstyle]
    #1
  \end{tcolorbox}
}

\newcommand{\warn}[1]{%
  \begin{tcolorbox}[warnstyle]
    #1
  \end{tcolorbox}
}

\newcommand{\exercise}[1]{%
  \begin{tcolorbox}[exercisestyle]
    #1
  \end{tcolorbox}
}

\newcommand{\defbox}[2]{%
  \begin{tcolorbox}[defstyle, title={#1}]
    #2
  \end{tcolorbox}
}

% -- Inline code ----------------------------------------------
\newcommand{\cmd}[1]{\texttt{\textcolor{kugreen!70!black}{#1}}}
\newcommand{\key}[1]{\fbox{\small\ttfamily #1}}

% -- Misc ----------------------------------------------------
\setlength{\parskip}{6pt}
\onehalfspacing

% ============================================================
\begin{document}

% -- Title page ----------------------------------------------
\begin{titlepage}
  \centering
  \vspace*{2cm}

  {\huge\bfseries\color{kugreen}
    Command Line \& HPC Cluster\\[0.4em]
    Tutorial for New Lab Members}

  \vspace{1.5cm}

  {\large\color{gray} From Zero to Running Your First Cluster Job}

  \vspace{2.5cm}

  \begin{tabular}{ll}
    \textbf{Lab}        & COMFHA Computational Lab \\[4pt]
    \textbf{Department} & Kasetsart University \\[4pt]
    \textbf{Audience}   & New lab members with no prior terminal experience \\[4pt]
    \textbf{Version}    & 1.0 (2026)
  \end{tabular}

  \vspace{2.5cm}

  \begin{tcolorbox}[enhanced, colback=kugreen!8, colframe=kugreen,
      left=12pt, right=12pt, top=8pt, bottom=8pt]
    \small
    This tutorial guides you step by step from opening a terminal for the
    first time to submitting your own job on the KU Cluster and Nontri HPC.
    No prior programming or Linux experience is required.
    Complete each section in order and attempt every exercise before moving on.
  \end{tcolorbox}

  \vfill
  {\small\color{gray} COMFHA Lab $\cdot$ Kasetsart University $\cdot$
   \url{comfha.ku.ac.th}}
\end{titlepage}

% -- TOC -----------------------------------------------------
\tableofcontents
\clearpage

% ============================================================
\part*{Part I \quad The Command Line}
\addcontentsline{toc}{part}{Part I: The Command Line}
% ============================================================

% -------------------------------------------------------------
\section{What Is a Terminal?}
% -------------------------------------------------------------

When you open a program like a file manager or a web browser, you interact with
it through icons and menus. A \textbf{terminal} (also called a \emph{shell},
\emph{command line}, or \emph{command prompt}) is a text-based way to interact
with your computer. Instead of clicking, you type instructions and the computer
responds.

\defbox{Why scientists use the terminal}{
  \begin{itemize}[noitemsep]
    \item Run programs that have no graphical interface (e.g.\ GROMACS, Python scripts)
    \item Automate repetitive tasks with a single command
    \item Connect to remote supercomputers (clusters)
    \item Work faster once you know a few commands
  \end{itemize}
}

\subsection{The Shell}

The program that reads your commands and runs them is called a \textbf{shell}.
The most common shell is \textbf{Bash} (Bourne Again SHell), which is the
default on most Linux systems and clusters. macOS uses \textbf{Zsh} (Z Shell)
by default, which is nearly identical for everyday use.

\subsection{Opening a Terminal}

\begin{center}
\begin{tabular}{>{\bfseries}lp{9cm}}
\toprule
Operating System & How to open a terminal \\
\midrule
macOS & \texttt{Cmd + Space} $\to$ type \texttt{Terminal} $\to$ press \texttt{Enter} \\
Ubuntu/Debian & \texttt{Ctrl + Alt + T} \\
Windows & Install \textit{Git Bash} or use \textit{Windows Subsystem for Linux (WSL)} \\
\bottomrule
\end{tabular}
\end{center}

When the terminal opens you will see something like this:

\begin{lstlisting}[style=terminal]
yourusername@computername ~ %
\end{lstlisting}

This is the \textbf{prompt}. It tells you who you are, which machine you are
on, and where you are in the file system. The \texttt{\%} (or \texttt{\$}) is
where you type your command.

\warn{
  Never close the terminal mid-task by just clicking the X. Type \cmd{exit}
  or press \key{Ctrl+D} to close it properly.
}

% -------------------------------------------------------------
\section{Your First Commands}
% -------------------------------------------------------------

Every command follows the same basic structure:

\begin{lstlisting}[style=code]
command  [options]  [arguments]
\end{lstlisting}

\begin{itemize}
  \item \textbf{command} --- the program you want to run
  \item \textbf{options} --- flags that change behavior, usually start with \texttt{-}
  \item \textbf{arguments} --- the target (file, directory, etc.)
\end{itemize}

\subsection{whoami, date, echo}

\begin{lstlisting}[style=bash]
whoami          # prints your username
date            # prints today's date and time
echo "Hello!"   # prints whatever you write
\end{lstlisting}

\begin{lstlisting}[style=terminal]
$ whoami
pao
$ date
Mon Jun 23 10:05:42 +07 2026
$ echo "Hello, lab!"
Hello, lab!
\end{lstlisting}

\tip{
  Text after a \texttt{\#} in a command is a \textbf{comment} --- the shell
  ignores it. Comments are used in this tutorial to explain what a line does.
}

\exercise{
  Open your terminal and run all three commands above. What is your username?
  What is today's date?
}

% -------------------------------------------------------------
\section{Navigating the File System}
% -------------------------------------------------------------

The file system is organized as a tree of folders (called \textbf{directories}).
On Linux and macOS the root of the tree is \texttt{/}. Your personal folder is
called the \textbf{home directory} and is written as \texttt{\textasciitilde}
(a tilde).

\begin{lstlisting}[style=bash]
/
+-- home/
|   \-- pao/          <- your home directory (~)
|       +-- Documents/
|       +-- Workspace/
|       \-- scripts/
+-- etc/
\-- usr/
\end{lstlisting}

\subsection{pwd --- Print Working Directory}

\cmd{pwd} tells you where you currently are.

\begin{lstlisting}[style=bash]
pwd
\end{lstlisting}
\begin{lstlisting}[style=terminal]
/Users/pao
\end{lstlisting}

\subsection{ls --- List Files}

\begin{lstlisting}[style=bash]
ls              # list files in current directory
ls -l           # long format (shows permissions, size, date)
ls -lh          # same but sizes in human-readable form (KB, MB)
ls -la          # include hidden files (names starting with .)
ls Documents/   # list a specific directory
\end{lstlisting}

\begin{lstlisting}[style=terminal]
$ ls -lh
total 48K
drwxr-xr-x  3 pao  staff   96B Jun 23 09:00 Documents
drwxr-xr-x  5 pao  staff  160B Jun 22 14:30 Workspace
-rw-r--r--  1 pao  staff  2.4K Jun 20 11:00 notes.txt
\end{lstlisting}

\subsection{cd --- Change Directory}

\begin{lstlisting}[style=bash]
cd Documents        # go into the Documents folder
cd ..               # go up one level (parent directory)
cd ../..            # go up two levels
cd ~                # go to your home directory (always works)
cd -                # go back to the previous directory
\end{lstlisting}

\tip{
  \textbf{Tab completion} is your best friend. Start typing a folder or file
  name and press \key{Tab} --- the shell will complete it for you. Press
  \key{Tab} twice to see all possible matches.
}

\exercise{
  \begin{enumerate}[noitemsep]
    \item Run \cmd{pwd} to see where you are.
    \item Run \cmd{ls -lh} to list your files.
    \item Navigate into one of the folders shown, then use \cmd{cd ..} to come back.
    \item Try typing \texttt{cd Doc} then press \key{Tab}. What happens?
  \end{enumerate}
}

% -------------------------------------------------------------
\section{Working with Files and Directories}
% -------------------------------------------------------------

\subsection{mkdir --- Make Directory}

\begin{lstlisting}[style=bash]
mkdir my_project             # create a folder
mkdir -p data/raw/outputs    # create nested folders at once
\end{lstlisting}

\subsection{touch --- Create an Empty File}

\begin{lstlisting}[style=bash]
touch notes.txt              # create an empty file called notes.txt
touch run.sh analysis.py     # create multiple files at once
\end{lstlisting}

\subsection{cp --- Copy}

\begin{lstlisting}[style=bash]
cp notes.txt notes_backup.txt        # copy a file
cp -r my_project/ my_project_old/    # copy a directory (-r = recursive)
\end{lstlisting}

\subsection{mv --- Move or Rename}

\begin{lstlisting}[style=bash]
mv notes.txt Documents/              # move file to Documents/
mv notes.txt summary.txt             # rename a file (same directory)
mv my_project/ Workspace/            # move a folder
\end{lstlisting}

\subsection{rm --- Remove (Delete)}

\warn{
  There is no Recycle Bin on the command line. Once you run \cmd{rm}, the
  file is gone permanently. Be very careful, especially with \cmd{rm -r}.
}

\begin{lstlisting}[style=bash]
rm notes.txt               # delete a file
rm -r old_data/            # delete a directory and everything inside it
rm -i file.txt             # ask for confirmation before deleting
\end{lstlisting}

\subsection{Quick Reference}

\begin{center}
\begin{tabular}{>{\ttfamily}lp{8cm}}
\toprule
\normalfont\textbf{Command} & \textbf{What it does} \\
\midrule
mkdir \textit{name}   & Create a directory \\
touch \textit{file}   & Create an empty file \\
cp \textit{a} \textit{b}    & Copy \textit{a} to \textit{b} \\
mv \textit{a} \textit{b}    & Move or rename \textit{a} to \textit{b} \\
rm \textit{file}      & Delete a file \\
rm -r \textit{dir}    & Delete a directory \\
\bottomrule
\end{tabular}
\end{center}

\exercise{
  \begin{enumerate}[noitemsep]
    \item Create a directory called \texttt{lab\_practice}.
    \item Inside it, create a file called \texttt{hello.txt}.
    \item Copy \texttt{hello.txt} to \texttt{hello\_copy.txt}.
    \item Rename \texttt{hello\_copy.txt} to \texttt{greeting.txt}.
    \item List the directory to confirm both files exist.
    \item Delete \texttt{greeting.txt}.
  \end{enumerate}
}

% -------------------------------------------------------------
\section{Viewing File Contents}
% -------------------------------------------------------------

\subsection{cat}

\cmd{cat} prints the entire file to the screen. Best for short files.

\begin{lstlisting}[style=bash]
cat notes.txt
\end{lstlisting}

\subsection{less}

\cmd{less} opens a file for scrolling --- press \key{Space} to go forward,
\key{b} to go back, \key{q} to quit. Best for long files.

\begin{lstlisting}[style=bash]
less trajectory_log.txt
\end{lstlisting}

\subsection{head and tail}

\begin{lstlisting}[style=bash]
head notes.txt         # show first 10 lines
head -n 20 notes.txt   # show first 20 lines
tail notes.txt         # show last 10 lines
tail -f output.log     # keep watching a file as it grows (useful for job logs)
\end{lstlisting}

\subsection{grep --- Search Inside Files}

\begin{lstlisting}[style=bash]
grep "error" output.log          # find lines containing "error"
grep -i "warning" output.log     # case-insensitive search
grep -n "Step" md.log            # show line numbers
\end{lstlisting}

\tip{
  \cmd{tail -f} is extremely useful when you have a simulation running --- it
  lets you watch the progress in real time without opening a new connection.
}

% -------------------------------------------------------------
\section{Editing Files with nano}
% -------------------------------------------------------------

\textbf{nano} is a simple text editor that runs inside the terminal.

\begin{lstlisting}[style=bash]
nano notes.txt          # open (or create) notes.txt for editing
\end{lstlisting}

Once inside nano, you can type normally. The keyboard shortcuts shown at the
bottom of the screen use \texttt{\^{}} to mean \key{Ctrl}:

\begin{center}
\begin{tabular}{lp{7cm}}
\toprule
\textbf{Shortcut} & \textbf{Action} \\
\midrule
\key{Ctrl+O} then \key{Enter} & Save (Write Out) \\
\key{Ctrl+X} & Exit nano \\
\key{Ctrl+K} & Cut the current line \\
\key{Ctrl+U} & Paste \\
\key{Ctrl+W} & Search \\
\key{Ctrl+\textbackslash} & Find and Replace \\
\bottomrule
\end{tabular}
\end{center}

\exercise{
  \begin{enumerate}[noitemsep]
    \item Open \texttt{nano hello.txt} inside your \texttt{lab\_practice} folder.
    \item Type: \texttt{My name is [your name] and I am learning the terminal.}
    \item Save with \key{Ctrl+O} then \key{Enter}.
    \item Exit with \key{Ctrl+X}.
    \item Verify with \cmd{cat hello.txt}.
  \end{enumerate}
}

% -------------------------------------------------------------
\section{Productivity Tips}
% -------------------------------------------------------------

\subsection{Command History}

Press the \key{$\uparrow$} arrow key to cycle through previous commands.
You do not need to retype them.

\begin{lstlisting}[style=bash]
history            # show all recent commands
history | grep cd  # find all times you used cd
\end{lstlisting}

\subsection{Ctrl Shortcuts}

\begin{center}
\begin{tabular}{ll}
\toprule
\textbf{Shortcut} & \textbf{What it does} \\
\midrule
\key{Ctrl+C}  & Cancel the currently running command \\
\key{Ctrl+Z}  & Pause the running command (resume with \texttt{fg}) \\
\key{Ctrl+L}  & Clear the screen \\
\key{Ctrl+A}  & Jump to beginning of the line \\
\key{Ctrl+E}  & Jump to end of the line \\
\key{Ctrl+D}  & Log out / close the terminal \\
\bottomrule
\end{tabular}
\end{center}

\subsection{The Pipe Operator \texttt{|}}

The pipe \texttt{|} sends the output of one command as input to another.
It is one of the most powerful features of the shell.

\begin{lstlisting}[style=bash]
ls -lh | grep ".txt"    # list only .txt files
cat data.log | tail -20 | grep "STEP"  # last 20 lines containing "STEP"
\end{lstlisting}

\subsection{Wildcards}

\begin{lstlisting}[style=bash]
ls *.txt          # list all files ending in .txt
rm output_*.log   # delete all files matching the pattern
cp *.py scripts/  # copy all Python files into scripts/
\end{lstlisting}

% ============================================================
\clearpage
\part*{Part II \quad SSH --- Connecting to Remote Computers}
\addcontentsline{toc}{part}{Part II: SSH --- Connecting to Remote Computers}
% ============================================================

% -------------------------------------------------------------
\section{What Is SSH?}
% -------------------------------------------------------------

\defbox{SSH (Secure Shell)}{
  SSH is a protocol that lets you log in to a remote computer over the
  network and use its terminal as if you were sitting in front of it.
  All communication is encrypted --- no one can intercept your commands
  or data.
}

In our lab you will use SSH to connect to two systems:

\begin{enumerate}
  \item \textbf{KU Cluster} (\texttt{ku-cluster}) --- general-purpose cluster
        at Kasetsart University. Used for MD simulations with GROMACS.
  \item \textbf{Nontri HPC} (\texttt{ku-ai}) --- GPU cluster (Nontri)
        at Kasetsart University. Used for GPU-accelerated jobs.
\end{enumerate}

\subsection{Basic SSH Syntax}

\begin{lstlisting}[style=bash]
ssh username@hostname
\end{lstlisting}

For example:

\begin{lstlisting}[style=bash]
ssh pao@158.108.25.40      # connect to KU Cluster
ssh pao@br1.paas.ku.ac.th  # connect to Nontri HPC
\end{lstlisting}

The first time you connect, SSH will ask you to verify the server's
fingerprint. Type \texttt{yes} and press \key{Enter}.

% -------------------------------------------------------------
\section{SSH Keys --- Password-Free Login}
% -------------------------------------------------------------

Typing a password every time is tedious and less secure. \textbf{SSH keys}
are a pair of files:

\begin{itemize}
  \item \textbf{Private key} --- stays on \emph{your} computer, never share it.
  \item \textbf{Public key} --- placed on the remote server. Anyone with the
        matching private key can log in.
\end{itemize}

\subsection{Step 1 --- Generate a Key Pair}

Run this on your \textbf{local machine} (your laptop or desktop):

\begin{lstlisting}[style=bash]
ssh-keygen -t ed25519 -C "your_email@example.com"
\end{lstlisting}

\begin{lstlisting}[style=terminal]
Generating public/private ed25519 key pair.
Enter file in which to save the key (/Users/pao/.ssh/id_ed25519):
Enter passphrase (empty for no passphrase):
Enter same passphrase again:
Your identification has been saved in /Users/pao/.ssh/id_ed25519
Your public key has been saved in /Users/pao/.ssh/id_ed25519.pub
\end{lstlisting}

\begin{itemize}
  \item Press \key{Enter} to accept the default file location.
  \item You may set a passphrase for extra security, or leave it empty.
\end{itemize}

This creates two files in \texttt{\textasciitilde/.ssh/}:

\begin{center}
\begin{tabular}{ll}
\toprule
\textbf{File} & \textbf{Description} \\
\midrule
\texttt{id\_ed25519}     & Your \textbf{private} key --- \textbf{never share this} \\
\texttt{id\_ed25519.pub} & Your \textbf{public} key --- this goes to the server \\
\bottomrule
\end{tabular}
\end{center}

\subsection{Step 2 --- Copy Your Public Key to the Server}

\begin{lstlisting}[style=bash]
# Easiest method (if ssh-copy-id is available):
ssh-copy-id -i ~/.ssh/id_ed25519.pub username@hostname

# Manual method:
cat ~/.ssh/id_ed25519.pub
# Copy the output, then on the remote server run:
# echo "PASTE_KEY_HERE" >> ~/.ssh/authorized_keys
\end{lstlisting}

After this you can log in without a password:

\begin{lstlisting}[style=bash]
ssh username@158.108.25.40   # no password prompt!
\end{lstlisting}

\warn{
  If you generate a \emph{separate} key pair for each cluster, give them
  different names during \cmd{ssh-keygen} (e.g.\ \texttt{id\_ed25519\_ku}).
  Use the \texttt{-i} flag with SSH to specify which key to use:\\[4pt]
  \texttt{ssh -i \textasciitilde/.ssh/id\_ed25519\_ku username@158.108.25.40}
}

\subsection{Step 3 --- Create an SSH Config File (Recommended)}

Instead of typing long commands, you can create a shortcut file at
\texttt{\textasciitilde/.ssh/config}:

\begin{lstlisting}[style=bash]
nano ~/.ssh/config
\end{lstlisting}

Add the following (replace \texttt{your\_username} with your actual username):

\begin{lstlisting}[style=code]
# KU Cluster
Host ku-cluster
    HostName 158.108.25.40
    User your_username
    IdentityFile ~/.ssh/id_ed25519_ku

# Nontri HPC
Host ku-ai
    HostName br1.paas.ku.ac.th
    User your_username
    IdentityFile ~/.ssh/id_ed25519
\end{lstlisting}

Save and exit nano (\key{Ctrl+O}, \key{Enter}, \key{Ctrl+X}).

Now you can connect with just:

\begin{lstlisting}[style=bash]
ssh ku-cluster      # instead of ssh -i ... username@158.108.25.40
ssh ku-ai           # instead of ssh -i ... username@br1.paas.ku.ac.th
\end{lstlisting}

% -------------------------------------------------------------
\section{Connecting to the KU Cluster}
% -------------------------------------------------------------

\subsection{Overview}

\begin{center}
\begin{tabular}{ll}
\toprule
\textbf{Property} & \textbf{Details} \\
\midrule
Hostname    & \texttt{158.108.25.40} \\
Short name  & \texttt{ku-cluster} (after SSH config is set up) \\
Scheduler   & SLURM \\
Recommended key & \texttt{id\_ed25519\_ku} \\
\bottomrule
\end{tabular}
\end{center}

\subsection{Logging In}

\begin{lstlisting}[style=bash]
ssh ku-cluster
\end{lstlisting}

You will see a welcome message followed by the cluster prompt:

\begin{lstlisting}[style=terminal]
Last login: Mon Jun 23 09:30:00 2026 from 10.31.74.100

[your_username@ku-cluster ~]$
\end{lstlisting}

\subsection{First Things to Do After Login}

\begin{lstlisting}[style=bash]
pwd                  # confirm your home directory path
ls                   # see what is already there
df -h ~              # check how much disk space you have
\end{lstlisting}

\subsection{Loading GROMACS}

Software on clusters is managed with the \cmd{module} system. You load the
software you need before running it:

\begin{lstlisting}[style=bash]
module avail                   # list all available software
module avail gromacs           # search for GROMACS specifically
module load gromacs/2020.4     # load a specific version
gmx --version                  # confirm it loaded
\end{lstlisting}

\tip{
  To avoid loading modules manually every time, add the \cmd{module load}
  line to your \texttt{\textasciitilde/.bashrc} file on the cluster.
}

\subsection{Logging Out}

\begin{lstlisting}[style=bash]
exit
\end{lstlisting}

% -------------------------------------------------------------
\section{Connecting to Nontri HPC (ku-ai)}
% -------------------------------------------------------------

\subsection{Overview}

\begin{center}
\begin{tabular}{ll}
\toprule
\textbf{Property} & \textbf{Details} \\
\midrule
Hostname     & \texttt{br1.paas.ku.ac.th} \\
Short name   & \texttt{ku-ai} (after SSH config is set up) \\
Scheduler    & SLURM \\
GPU queue    & \texttt{gpuq} \\
GROMACS      & 2022.6, 2024.1 \\
\bottomrule
\end{tabular}
\end{center}

\subsection{Logging In}

\begin{lstlisting}[style=bash]
ssh ku-ai
\end{lstlisting}

\begin{lstlisting}[style=terminal]
[your_username@nontri ~]$
\end{lstlisting}

\subsection{Loading GROMACS on Nontri}

\begin{lstlisting}[style=bash]
module avail gromacs
module load gromacs/2024.1    # GPU-enabled version
gmx --version
\end{lstlisting}

\warn{
  Nontri uses GPU nodes for GROMACS. Always request a \texttt{gpuq} partition
  in your SLURM script (see Section~\ref{sec:slurm}). Running a
  GPU-optimized binary on CPU-only nodes will either fail or run very slowly.
}

% -------------------------------------------------------------
\section{Transferring Files Between Your Computer and a Cluster}
% -------------------------------------------------------------

\subsection{scp --- Secure Copy}

\cmd{scp} works like \cmd{cp} but over SSH.

\begin{lstlisting}[style=bash]
# Upload: local -> cluster
scp my_script.py ku-cluster:~/Workspace/

# Download: cluster -> local
scp ku-cluster:~/outputs/result.xvg ./

# Copy a whole directory (recursive)
scp -r local_folder/ ku-cluster:~/Workspace/

# Same for Nontri
scp analysis.py ku-ai:~/project/
\end{lstlisting}

\subsection{rsync --- Smart Sync (Recommended)}

\cmd{rsync} only transfers files that have changed, making it much faster
for large directories.

\begin{lstlisting}[style=bash]
# Sync a local folder to the cluster
rsync -avz --progress local_folder/ ku-cluster:~/Workspace/local_folder/

# Download results back to your machine
rsync -avz ku-cluster:~/outputs/ ./outputs/
\end{lstlisting}

\begin{center}
\begin{tabular}{ll}
\toprule
\textbf{Flag} & \textbf{Meaning} \\
\midrule
\texttt{-a} & Archive mode (preserves timestamps, permissions, etc.) \\
\texttt{-v} & Verbose (show what is being transferred) \\
\texttt{-z} & Compress data during transfer (faster on slow networks) \\
\texttt{--progress} & Show transfer progress \\
\bottomrule
\end{tabular}
\end{center}

\tip{
  Always use \cmd{rsync} instead of \cmd{scp} when syncing simulation output
  folders --- trajectory files are large and rsync avoids re-sending what
  already arrived.
}

% ============================================================
\clearpage
\part*{Part III \quad Running Jobs on the Cluster (SLURM)}
\addcontentsline{toc}{part}{Part III: Running Jobs on the Cluster}
\label{sec:slurm}
% ============================================================

% -------------------------------------------------------------
\section{What Is a Job Scheduler?}
% -------------------------------------------------------------

A cluster is a shared resource --- many researchers use it at the same time.
A \textbf{scheduler} (our clusters use \textbf{SLURM}) manages the queue of
jobs and decides which job runs on which node and when.

You submit a \textbf{job script} --- a text file that describes what resources
you need and what commands to run. SLURM puts your job in the queue and runs
it when resources are available.

\defbox{Key concepts}{
  \begin{description}[noitemsep]
    \item[Node] One physical computer in the cluster.
    \item[Core / CPU] A single processing unit on a node.
    \item[Partition] A group of nodes (e.g.\ GPU nodes, high-memory nodes).
    \item[Job] Your submitted work unit.
    \item[Queue] The list of waiting jobs.
    \item[Walltime] The maximum time your job is allowed to run.
  \end{description}
}

% -------------------------------------------------------------
\section{Checking the Cluster Status}
% -------------------------------------------------------------

\begin{lstlisting}[style=bash]
sinfo                        # see all partitions and node status
squeue                       # see all jobs currently in the queue
squeue -u your_username      # see only YOUR jobs
squeue --start               # see estimated start times
\end{lstlisting}

\begin{lstlisting}[style=terminal]
$ squeue -u pao
JOBID  PARTITION  NAME       USER  ST  TIME        NODES NODELIST
12345  main       run_nvt    pao   R   0:05:23     1     node04
12346  gpuq       run_gpu    pao   PD  0:00:00     1     (Priority)
\end{lstlisting}

Common \textbf{ST} (state) values:

\begin{center}
\begin{tabular}{ll}
\toprule
\textbf{State} & \textbf{Meaning} \\
\midrule
\texttt{R}  & Running \\
\texttt{PD} & Pending (waiting in queue) \\
\texttt{CG} & Completing \\
\texttt{F}  & Failed \\
\bottomrule
\end{tabular}
\end{center}

% -------------------------------------------------------------
\section{Writing a SLURM Job Script}
% -------------------------------------------------------------

A job script is a regular shell script with special comments at the top
that start with \texttt{\#SBATCH}. These tell SLURM what resources you need.

\subsection{Template for KU Cluster (CPU job)}

\begin{lstlisting}[style=bash]
#!/bin/bash
#SBATCH --job-name=my_simulation     # name shown in squeue
#SBATCH --output=slurm_%j.out        # stdout log (%j = job ID)
#SBATCH --error=slurm_%j.err         # stderr log
#SBATCH --nodes=1                    # number of nodes
#SBATCH --ntasks=1                   # number of MPI tasks
#SBATCH --cpus-per-task=6            # CPU cores per task
#SBATCH --mem=8G                     # total memory
#SBATCH --time=24:00:00              # max walltime (HH:MM:SS)
#SBATCH --partition=main             # partition name

# -- Environment setup --------------------------------------
module load gromacs/2020.4

# -- Change to the directory you submitted from -------------
cd $SLURM_SUBMIT_DIR

# -- Your commands ------------------------------------------
echo "Job started on $(hostname) at $(date)"

gmx mdrun -v -deffnm run -ntmpi 1 -ntomp 6

echo "Job finished at $(date)"
\end{lstlisting}

Save this file as, for example, \texttt{submit.sh}.

\subsection{Template for Nontri HPC (GPU job)}

\begin{lstlisting}[style=bash]
#!/bin/bash
#SBATCH --job-name=gpu_simulation
#SBATCH --output=slurm_%j.out
#SBATCH --error=slurm_%j.err
#SBATCH --nodes=1
#SBATCH --ntasks=1
#SBATCH --cpus-per-task=4
#SBATCH --gres=gpu:1                 # request 1 GPU
#SBATCH --mem=16G
#SBATCH --time=48:00:00
#SBATCH --partition=gpuq             # GPU partition on Nontri

module load gromacs/2024.1

cd $SLURM_SUBMIT_DIR

echo "Job started on $(hostname) at $(date)"

# GPU-enabled GROMACS run
gmx mdrun -v -deffnm run -ntmpi 1 -ntomp 4 -gpu_id 0

echo "Job finished at $(date)"
\end{lstlisting}

\subsection{Common \#SBATCH Options}

\begin{center}
\begin{tabular}{>{\ttfamily}lp{8cm}}
\toprule
\normalfont\textbf{Option} & \textbf{Description} \\
\midrule
--job-name=\textit{name}       & Job name (appears in squeue) \\
--output=\textit{file}         & Standard output log file \\
--error=\textit{file}          & Error log file \\
--nodes=\textit{N}             & Number of nodes \\
--ntasks=\textit{N}            & Number of parallel tasks (MPI) \\
--cpus-per-task=\textit{N}     & CPU cores per task (OpenMP) \\
--mem=\textit{size}            & Total memory (e.g.\ 8G) \\
--time=\textit{HH:MM:SS}       & Maximum runtime \\
--partition=\textit{name}      & Which partition to use \\
--gres=gpu:\textit{N}          & Number of GPUs to request \\
\bottomrule
\end{tabular}
\end{center}

% -------------------------------------------------------------
\section{Submitting and Managing Jobs}
% -------------------------------------------------------------

\begin{lstlisting}[style=bash]
# Submit your job script
sbatch submit.sh

# Check the status of your jobs
squeue -u your_username

# Cancel a job (use the JOBID from squeue)
scancel 12345

# Cancel ALL your jobs
scancel -u your_username

# See details and resource usage of a completed job
sacct -j 12345 --format=JobID,JobName,State,Elapsed,MaxRSS

# See estimated start time
squeue --start -u your_username
\end{lstlisting}

\begin{lstlisting}[style=terminal]
$ sbatch submit.sh
Submitted batch job 12345

$ squeue -u pao
JOBID  PARTITION  NAME           USER  ST  TIME      NODES
12345  main       my_simulation  pao   R   0:01:45   1
\end{lstlisting}

\subsection{Monitoring Your Running Job}

Your job writes output to the log files you specified in the script.
Watch them in real time with:

\begin{lstlisting}[style=bash]
tail -f slurm_12345.out      # replace 12345 with your actual job ID
\end{lstlisting}

Press \key{Ctrl+C} to stop watching the file.

\exercise{
  \begin{enumerate}[noitemsep]
    \item Log in to one of the clusters.
    \item Create a file called \texttt{test\_job.sh} using nano.
    \item Paste a simple job script that runs \texttt{echo "Hello from the cluster!"}.
    \item Submit it with \cmd{sbatch test\_job.sh}.
    \item Watch the queue with \cmd{squeue -u your\_username}.
    \item When the job finishes, read the output file with \cmd{cat}.
  \end{enumerate}
}

% -------------------------------------------------------------
\section{Keeping Jobs Running After You Log Out}
% -------------------------------------------------------------

Jobs submitted with \cmd{sbatch} continue running even after you disconnect ---
this is the whole point of a scheduler. However, if you run a program
\emph{directly} in the terminal (not via sbatch) and you disconnect,
it will be killed.

\begin{tcolorbox}[defstyle, title={Use screen or tmux for interactive sessions}]
  If you need an interactive session that survives disconnection, use
  \texttt{screen} or \texttt{tmux}:
\begin{lstlisting}[style=bash]
screen -S mysession     # start a named screen session
# ... run your commands ...
# Detach: Ctrl+A then D
screen -r mysession     # reattach after reconnecting
\end{lstlisting}
  However, for all production jobs in our lab, always use \texttt{sbatch}.
\end{tcolorbox}

% ============================================================
\clearpage
\part*{Appendix \quad Quick Reference Card}
\addcontentsline{toc}{part}{Appendix: Quick Reference Card}
% ============================================================

\section*{A.1 --- Essential Commands}

\begin{center}
\renewcommand{\arraystretch}{1.3}
\begin{tabular}{>{\ttfamily}lp{9cm}}
\toprule
\normalfont\textbf{Command} & \textbf{What it does} \\
\midrule
pwd              & Print current directory \\
ls -lh           & List files (human-readable sizes) \\
cd \textit{dir}        & Change to directory \\
cd ..            & Go up one level \\
mkdir \textit{name}    & Create directory \\
touch \textit{file}    & Create empty file \\
cp \textit{a} \textit{b}     & Copy a to b \\
mv \textit{a} \textit{b}     & Move/rename a to b \\
rm \textit{file}       & Delete file \\
cat \textit{file}      & Print file contents \\
less \textit{file}     & Scroll through file (q to quit) \\
tail -f \textit{file}  & Watch file in real time \\
grep \textit{pat} \textit{file} & Search for pattern in file \\
nano \textit{file}     & Open text editor \\
history          & Show command history \\
\bottomrule
\end{tabular}
\end{center}

\section*{A.2 --- SSH \& File Transfer}

\begin{center}
\renewcommand{\arraystretch}{1.3}
\begin{tabular}{>{\ttfamily}lp{8cm}}
\toprule
\normalfont\textbf{Command} & \textbf{What it does} \\
\midrule
ssh ku-cluster                        & Connect to KU Cluster \\
ssh ku-ai                             & Connect to Nontri HPC \\
ssh-keygen -t ed25519                 & Generate SSH key pair \\
ssh-copy-id user@host                 & Install public key on server \\
scp file.txt ku-cluster:\textasciitilde/        & Upload file to cluster \\
scp ku-cluster:\textasciitilde/file.txt .       & Download file from cluster \\
rsync -avz dir/ ku-cluster:\textasciitilde/dir/ & Sync directory to cluster \\
exit                                  & Log out \\
\bottomrule
\end{tabular}
\end{center}

\section*{A.3 --- SLURM Job Management}

\begin{center}
\renewcommand{\arraystretch}{1.3}
\begin{tabular}{>{\ttfamily}lp{8cm}}
\toprule
\normalfont\textbf{Command} & \textbf{What it does} \\
\midrule
sbatch submit.sh              & Submit a job script \\
squeue -u \textit{username}          & Check your jobs \\
squeue --start                & See estimated start times \\
scancel \textit{JOBID}               & Cancel a specific job \\
scancel -u \textit{username}         & Cancel all your jobs \\
sinfo                         & Show partition/node status \\
sacct -j \textit{JOBID}              & Show finished job details \\
module avail                  & List available software \\
module load \textit{name/version}    & Load software \\
module list                   & List currently loaded modules \\
\bottomrule
\end{tabular}
\end{center}

\section*{A.4 --- Cluster Connection Details}

\begin{center}
\renewcommand{\arraystretch}{1.4}
\begin{tabular}{llll}
\toprule
\textbf{Cluster} & \textbf{SSH alias} & \textbf{Hostname} & \textbf{Scheduler} \\
\midrule
KU Cluster   & \texttt{ku-cluster} & \texttt{158.108.25.40}     & SLURM \\
Nontri HPC   & \texttt{ku-ai}      & \texttt{br1.paas.ku.ac.th} & SLURM (\texttt{gpuq}) \\
\bottomrule
\end{tabular}
\end{center}

\vspace{2em}
\begin{tcolorbox}[enhanced, colback=kugreen!8, colframe=kugreen,
    title={\bfseries Getting Help}]
  \begin{itemize}[noitemsep]
    \item \textbf{Manual pages:} type \texttt{man command} (e.g.\ \texttt{man ls}) for full documentation
    \item \textbf{Quick help:} type \texttt{command --help} for a short summary
    \item \textbf{SLURM docs:} \url{https://slurm.schedmd.com/documentation.html}
    \item \textbf{GROMACS docs:} \url{https://manual.gromacs.org}
    \item \textbf{Ask your supervisor or senior lab members} --- we are here to help!
  \end{itemize}
\end{tcolorbox}

\end{document}
